%\begin{name}
%	{Biên soạn: Thầy Trần Chiến \& Phản biện: Thầy Đỗ Minh Phúc}
%	{Đề chuẩn hóa}
%\end{name}

\caulc
\Opensolutionfile{ans}[ans/ans-cau-TN-1]

\begin{ex}%[EX chuẩn hóa]%[Trần Chiến]%[1D7Y2-1]
	Đạo hàm của hàm số $y = x\sin x$ bằng
		\choice
		{\True $y' = \sin x + x\cos x$}
		{$y' = \sin x - x\cos x$}
		{$y' = x\cos x$}
		{$y' =  - x\cos x$}
		\loigiai{
			Ta có $y'=\sin x+x\cos x.$
		}
	\end{ex}

\begin{ex}%[EX chuẩn hóa]%[Trần Chiến]%[1D7Y2-1]
	Hàm số $y = \sin \left( \dfrac{\pi }{2} - 2x \right)$ có đạo hàm bằng
	\choice
	{$y' = 2\sin 2x$}
	{\True $y' =  - 2\sin 2x$}
	{$y' = \cos \left( \dfrac{\pi }{2} - 2x \right)$}
	{$y' = 2\cos \left( \dfrac{\pi }{2} - 2x \right)$}
	\loigiai{
		Ta có $y=\cos 2x$ nên $y'=-2\sin 2x$.
	}
\end{ex}



\begin{ex}%[EX chuẩn hóa]%[Trần Chiến]%[1D7Y1-1]
	Chọn khẳng định đúng trong các khẳng định sau
	\choice
	{Hàm số $y=f(x)$ liên tục tại điểm $x_0$ thì có đạo hàm tại điểm đó}
	{\True Hàm số $y=f(x)$ có đạo hàm tại điểm $x_0$ thì liên tục tại điểm đó}
	{Hàm số $y=f(x)$ xác định tại điểm $x_0$ thì có đạo hàm tại điểm đó}
	{Hàm số $y=f(x)$ luôn có đạo hàm tại mọi điểm thuộc tập xác định của nó}
	\loigiai{Nếu hàm số $y=f(x)$ có đạo hàm tại điểm $x_0$ thì liên tục tại điểm đó.
	}
\end{ex}

\begin{ex}%[EX chuẩn hóa]%[Trần Chiến]%[1D7Y1-2]
	Số gia $\Delta y$ của hàm số $y = x^2 + 2x - 5$ tại điểm $x_0 = 1$ là
	\choice
	{$(\Delta x)^2 - 4\Delta x$}
	{\True $\left( \Delta x \right)^2 + 4\Delta x$}
	{$\left(\Delta x \right)^2 - 2\Delta x$}
	{$\left( \Delta x \right)^2 + 2\Delta x - 5$}
	\loigiai{
		Ta có $$\begin{aligned} \Delta y&=f(1+\Delta x)-f(1)=\left(1+\Delta x \right)^2+2(1+\Delta x)-5+2 \\&= \left(\Delta x \right)^2+4\Delta x.  \end{aligned} $$
	}
\end{ex}

\begin{ex}%[EX chuẩn hóa]%[Trần Chiến]%[1D7Y2-1]
	Đạo hàm của hàm số $y = \dfrac{x^2 - 2x - 1}{x - 2}$ bằng 
	\choice
	{\True $y' = \dfrac{x^2 - 4x + 5}{(x - 2)^2}$}
	{$y' = \dfrac{x^2 - 6x + 5}{(x - 2)^2}$}
	{$y' = \dfrac{x^2 - 6x + 4}{(x - 2)^2}$}
	{$y' = \dfrac{x^2 - 6x - 1}{(x - 2)^2}$}
	\loigiai{
		Ta có 
		$$y'=\dfrac{(2x-2)(x-2)-(x^2-2x-1)}{(x-2)^2}=\dfrac{x^2-4x+5}{(x-2)^2}.$$
	}
\end{ex}

\begin{ex}%[EX chuẩn hóa]%[Trần Chiến]%[1D7B2-1]
	Đạo hàm cấp hai của hàm số $y = \tan x$ là
	\choice
	{$y'' = 2\tan x \left(1 - \tan^2 x \right)$}
	{\True $y'' = 2\tan x \left(1 + \tan^2 x \right)$}
	{$y'' =  - 2\tan x \left(1 - \tan^2 x \right)$}
	{$y'' =  - 2\tan \left(1 + \tan^2 x \right)$}
	\loigiai{
		$y'=\tan^2 x + 1$;\hspace{1cm} $y'' = 2\tan x \cdot (\tan x)' = 2\tan x \cdot (\tan^2 x + 1)$.
	}
\end{ex} 
\begin{ex}%[EX chuẩn hóa]%[Trần Chiến]%[1D7B3-2]
	Cho hàm số $y = \dfrac{1}{{1 - x}}$. Chọn khẳng định đúng trong các khẳng định sau:
	\choice
	{$y'' + y^3 = 0$}
	{$y'' - y^3 = 0$}
	{$y'' + 2y^3 = 0$}
	{\True $y'' - 2y^3 = 0$}
	\loigiai{
		Ta có $y'=\dfrac{1}{(1-x)^2},\ y''=\dfrac{2}{(1-x)^3}$. Suy ra $y''-2y^3=0$.
	}
\end{ex}


\begin{ex}%[EX chuẩn hóa]%[Trần Chiến]%[1H8Y7-1]
	Cho hình chóp tam giác đều $S.ABC$. Khẳng định nào sau đây là \textbf{sai} về hình chóp đã cho? 
	\choice
	{Tam giác $ABC$ là tam giác đều}
	{Các cạnh bên hợp với đáy các góc bằng nhau}
	{\True Các mặt bên là các tam giác đều}
	{Các mặt bên hợp với đáy các góc bằng nhau}
	\loigiai{
		Hình chóp tam giác đều có các mặt bên là các tam giác cân.
	}
\end{ex}


\begin{ex}%[EX chuẩn hóa]%[Trần Chiến]%[1D7K2-4]
	Cho hàm số $y = 4x + 2\cos 2x$ có đồ thị là $(C)$. Hoành độ của các điểm trên $(C)$ mà tại đó tiếp tuyến của $(C)$ song song hoặc trùng với trục hoành là
	\choice
	{\True $x = \dfrac{\pi}{4} + k\pi \enspace \left( k \in \mathbb{Z} \right)$}
	{$x = \pi  + k\pi \enspace \left( k \in \mathbb{Z} \right)$}
	{$x = \dfrac{\pi}{2} + k\pi \enspace \left( k \in \mathbb{Z} \right)$}
	{$x = k2\pi \enspace \left( k \in \mathbb{Z} \right)$}
	\loigiai{
		$y'=4 - 4\sin 2x$. \\
		Gọi $M(x_0; y_0)$ là tọa độ tiếp điểm của $(C)$ và tiếp tuyến cần tìm.\\
		Trục hoành có phương trình $y = 0$. Để tiếp tuyến của $(C)$ song song hoặc trùng với trục hoành thì hệ số góc của tiếp tuyến cần tìm là $y'(x_0) = 0$. Suy ra \\
		$4 - 4\sin 2x_0 = 0 \Leftrightarrow \sin 2x_0 = 1 \Leftrightarrow 2x_0 = \dfrac{\pi}{2} + k2\pi \Leftrightarrow x_0 = \dfrac{\pi}{4} + k\pi \ (k\in\mathbb{Z})$.
	}
\end{ex}

\begin{ex}%[EX chuẩn hóa]%[Trần Chiến]%[1D7K1-1]
	Cho hàm số $ y=f(x)= \begin{cases} \sqrt{4x+1} & \text{khi } x>0 \\ x+1 & \text{khi } x \leq 0 \end{cases}$. Khẳng định nào là đúng về đạo hàm của hàm số $f(x)$ tại $x=0$?
	\choice
	{$f'(0) =  - 1$}
	{$f'(0) = 0$}
	{$f'(0)  = 1$}
	{\True Không tồn tại}
	\loigiai{
		$\displaystyle \lim_{x \rightarrow 0^+} \dfrac{f(x) - f(0)}{x-0} 
		= \lim_{x \rightarrow 0^+} \dfrac{\sqrt{4x+1} - 1}{x} 
		=\lim_{x \rightarrow 0^+} \dfrac{4x + 1 - 1}{x \left(\sqrt{4x+1} + 1\right)} = 2$;\\
		$\displaystyle \lim_{x \rightarrow 0^-} \dfrac{f(x) - f(0)}{x-0} 
		= \lim_{x \rightarrow 0^-} \dfrac{x+1 - 1}{x-0} = 1$.\\
		Do $\displaystyle \lim_{x \rightarrow 0^+}  \dfrac{f(x) - f(0)}{x-0} \ne \lim_{x \rightarrow 0^-} \dfrac{f(x) - f(0)}{x-0}$ nên $\displaystyle \lim_{x \rightarrow 0}  \dfrac{f(x) - f(0)}{x-0}$ không tồn tại.
	}
\end{ex}

\begin{ex}%[EX chuẩn hóa]%[Trần Chiến]%[1H8Y4-1]
	Cho hình chóp tam giác đều $S.ABC$. Khẳng định nào sau đây là \textbf{sai} về hình chóp đã cho? 
	\choice
	{Tam giác $ABC$ là tam giác đều}
	{Các cạnh bên hợp với đáy các góc bằng nhau}
	{\True Các mặt bên là các tam giác đều}
	{Các mặt bên hợp với đáy các góc bằng nhau}
	\loigiai{
		Hình chóp tam giác đều có các mặt bên là các tam giác cân.
	}
\end{ex}
\begin{ex}%[EX chuẩn hóa]%[Trần Chiến]%[1H8B2-2]
	Cho hình lăng trụ đứng $ABCD.A'B'C'D'$ có đáy $ABCD$ là hình vuông. Khẳng định nào sau đây là đúng?
	\choice
	{$A'C \perp (B'BD)$}
	{$A'C \perp (B'C'D)$}
	{\True $AC \perp (B'BD')$}
	{$AC \perp (B'CD')$}
	\loigiai{\immini{
			$\begin{cases} AC \perp BD \\ AC \perp BB' \end{cases} \Rightarrow AC \perp (BB'D'D)$ hay $AC \perp (BB'D')$.
		}{	
			\begin{tikzpicture}[declare function={r=3;},font=\footnotesize,scale=0.8, line join=round, line cap=round, >=stealth]
			\path (0:0) coordinate (A)
			(0:r) coordinate (B)
			++(37:{0.65*r}) coordinate (C)
			++(180:r) coordinate (D)
			($(B)!1/2!(D)$) coordinate (M)
			
			\foreach \x in {A,B,C,D}{(\x)++(90:r) coordinate (\x')};
			\draw[dash pattern=on 2pt off 2 pt] (D')--(D)--(A) (D)--(C) (A)--(C) (B)--(D);
			\draw (A)--(B)--(C) (A)--(A') (B)--(B') (C)--(C') (B')--(D') (A')--(B')--(C')--(D')--cycle;
			\foreach \t/\g in {A/180,B/0,C/0,D/-80,A'/180,B'/90,C'/0,D'/180}{
				\draw[fill=black] (\t) circle (1pt) node[shift={(\g:7pt)},font=\scriptsize]{$ \t $};
			}
		\foreach \x/\y/\z in {B/M/C} \draw pic[draw,angle radius=2.5mm] {right angle= \x--\y--\z};  
			\end{tikzpicture}
	}}
\end{ex}


\Closesolutionfile{ans}
\indapan{7}{ans/ans-cau-TN-1}

\cauds
\Opensolutionfile{ans}[ans/ans-cau-DS-1]


\begin{ex}%[EX chuẩn hóa]%[Trần Chiến]%[1H8B4-1]
	Xét các khẳng định nào sau
	\choiceTF
	{\True $\heva{&(\alpha) \cap (\beta) = d\\& (\alpha) \perp (P) \\& (\beta) \perp (P)}\Rightarrow d \perp (P)$}
	{\True $\heva{&(\alpha)\parallel (\beta) \\& ( P ) \perp ( \alpha  )}
		\Rightarrow ( P ) \perp ( \beta  )$}
	{$\heva{&(\alpha) \perp (\beta) \\& a \subset (\alpha) \\& b \subset ( \beta  )} \Rightarrow a \perp b$}
	{$\heva{&(\alpha) \perp (\beta) \\& a \subset (\alpha)} \Rightarrow a \perp (\beta)$}
	\loigiai{
		\begin{center}
			\begin{tikzpicture}[scale=0.7,font=\footnotesize,line join=round,line cap=round,>=stealth]
				\path
				(0,0) coordinate (B)
				(1,1.4) coordinate (A)
				(4,0) coordinate (C)
				($(C)-(B)+(A)$) coordinate (D)	
				($(A)+(90:3.5)$) coordinate (A')
				($(B)-(A)+(A')$) coordinate (B')
				($(C)-(A)+(A')$) coordinate (C')
				($(D)-(A)+(A')$) coordinate (D')
				;
				\draw (B')--(B)--(C)--(D)--(D')--(A')--(B')--(C')--(D')--cycle (C')--(C) (B')--(C);
				\draw[dashed] (A')--(A)--(B) (A)--(D) ;
				\foreach \p/\q in {A/160,B/-135,C/-45,D/0,A'/90,B'/180,C'/-20,D'/90}
				\fill[black] (\p)node[shift={(\q:3mm)}]{$\p$} circle (1.0pt);	
			\end{tikzpicture}
		\end{center}
		\begin{itemchoice}
			\itemch Đúng, vì hai mặt phẳng cùng vuông góc với mặt phẳng thứ 3 thì giao tuyến của chúng cũng vuông góc.
			\itemch Đúng, vì một mặt phẳng vuông góc với một trong hai mặt phẳng song song thì sẽ vuông góc với mặt phẳng còn lại.
			\itemch Sai vì, xét hỉnh lập phương $ABCD.A'B'C'D'$ có
			\[\heva{& (A'B'C'D') \perp (BCC'B') \\& B'D' \subset (A'B'C'D'), B'C \subset (BCC'B')  \\ & BC', B'D' \text{ không vuông góc}}\]
			\itemch Sai vì, xét hỉnh lập phương $ABCD.A'B'C'D'$ có
			\[\heva{& (A'B'C'D') \perp (BCC'B') \\& B'D' \subset (A'B'C'D')  \\ &  B'D' \text{ không vuông góc với } (BCC'B')}\]
		\end{itemchoice}
	}
\end{ex}

\begin{ex}%[EX chuẩn hóa]%[Trần Chiến]%[1D7B3-2]
	Cho hàm số $y = \sin ^2 x$. 
	\choiceTF
	{\True $y'= 2\sin x \cos x$}
	{$y''=-2\cos 2x$}
	{$(y')^2 + (y'')^2 = 1$}
	{\True $(y')^2 + (1 - 2y)^2 = 1$}
	\loigiai{Ta có
		\begin{itemchoice}
			\itemch Đúng, vì $y' = 2\sin x \cdot (\sin x)' = 2\sin x \cos x = \sin 2x$.
			\itemch Sai, vì $y''=2\cos 2x$.
			\itemch Sai, vì $(y')^2 + (y'')^2 = \sin^2 2x + 4\cos^2 2x = 1+3\cos^2 2x \ne 1$.
			\itemch Đúng, vì $(y')^2 + (1-2y)^2 = \sin^2 2x + (1-2\sin^2 x)^2 = \sin^2 2x + \cos^2 2x = 1$.
		\end{itemchoice}
	}
\end{ex}

\begin{ex}%[EX chuẩn hóa]%[Trần Chiến]%[1D7B2-1]
	Cho hai hàm số $f(x)= x + 2$ và $g(x) = x^2 - 2x + 3$.
	\choiceTF
	{\True $f'(x)=1$}
	{$g'(x)=x-2$}
	{$[f(g(1))]'=2$}
	{\True $[g(f(1))]'= 4$}
	\loigiai{
		Ta có 
		\begin{itemchoice}
			\itemch $f'(x)=1$.
			\itemch $g'(x)=2x-2$.
			\itemch $[f(g(x))]'= f'(g(x)) \cdot g'(x) = 1\cdot (2x-2) =2x-2$, suy ra $[f(g(0))]'=-2$
			\itemch $[g(f(x))]'=g'\left(f(x)\right) \cdot f'(x) = \left[2(x+2)-2\right]\cdot 1= 2x+2$, suy ra $[g(f(1))]' = 4$
		\end{itemchoice}
	}
\end{ex}

\begin{ex}%[EX chuẩn hóa]%[Trần Chiến]%[1H8K4-3]
	Cho tứ diện $S.ABC$ có các tam giác $SAB, SAC$ và $ABC$ vuông cân tại $A$, $SA = a$. Gọi $I$ là trung điểm của $BC$ và $ \alpha $ là góc giữa hai mặt phẳng $(SBC)$ và $(ABC)$.
	\choiceTF
	{\True $SA \perp BC$}
	{$AI=a\sqrt{2}$}
	{$AI \perp (SBC)$}
	{\True $\tan \alpha  = \sqrt 2 $}
	\loigiai{
		\begin{center}
			\begin{tikzpicture}[scale=1, font=\footnotesize, line join=round, line cap=round, >=stealth]
		\path
		(0,0) coordinate (A)
		(3,0) coordinate (C)
		(1,-1) coordinate (B)
		($(B)!0.5!(C)$) coordinate (I)
		%(intersection of A--M and C--N) coordinate (H)
		($(A)+(0,3)$) coordinate (S);
		\draw (B)--(S)--(A)--(B)--(C)--(S)--(I); \draw[dashed] (C)--(A)--(I);
		\foreach \i/\g in {S/90, A/-135, B/-90, C/-20, I/-45}{
			\draw[fill=black] (\i) circle (1.2pt) +(\g:0.25) node{$\i$};}
		
		\draw pic[draw,angle eccentricity=1.7, angle radius=0.4cm]{angle=S--I--A};
		\draw pic[draw,angle radius=2mm]{right angle=A--I--B};
		\draw pic[draw,angle radius=2mm]{right angle=S--I--C};
		\end{tikzpicture}
		\end{center}
		\begin{itemchoice}
			\itemch Ta có 
			$$\heva{ &SA \perp AB \\& SA \perp AC } \Rightarrow SA \perp (ABC) \Rightarrow SA \perp BC. (1)$$
			\itemch Ta có 	$AB = AC = SA = a$; $BC = AB \sqrt{2} = a\sqrt{2}$; $AI = \dfrac{BC}{2} = \dfrac{a\sqrt{2}}{2}$.
			\itemch Sai
			\itemch Do tam giác $ABC$ cân tại $A$ nên $AI \perp BC$. $(2)$\\
			Từ $(1)$ và $(2)$ suy ra $BC \perp (SAI) \Rightarrow BC \perp SI$.\\
			Ta có $\heva{&(SBC) \cap (ABC) = BC \\& SI \subset (SBC), SI \perp BC \\& AI \subset (ABC), AI \perp BC }\Rightarrow \left[ (SBC), (ABC) \right] = (SI, AI) = \widehat{SIA}$.\\
			$\tan \alpha = \tan \widehat{SIA} = \dfrac{SA}{AI} = a \cdot \dfrac{2}{a\sqrt{2}} = \sqrt{2}$.
		\end{itemchoice}
	}
\end{ex}

\Closesolutionfile{ans}
\indapan{3}{ans/ans-cau-DS-1}
\caukq
\Opensolutionfile{ans}[ans/ans-cau-KQ-1]

\begin{ex}%[EX chuẩn hóa]%[Trần Chiến]%[1D7Y2-2]
	Tiếp tuyến với đồ thị của hàm số $f(x) = \dfrac{3}{2x - 1}$ tại điểm có hoành độ $x_0 = 2$ có hệ số góc là bao nhiêu (làm tròn kết quả đến hàng phần chục)?
	\shortans[]{$-0{,} 7$}
	\loigiai{
		Ta có $f'(x)=-\dfrac{6}{(2x-1)^2}$ nên $f'(2)=-\dfrac{6}{9}=-\dfrac{2}{3}\approx -0{,} 7$.
	}
\end{ex}


\begin{ex}%[EX chuẩn hóa]%[Trần Chiến]%[1D6K4-4]
	Tính tổng tất cả các nghiệm của phương trình $ \log_4x^2+\log_2(5-x)=\log_2(x+3) $ (làm tròn kết quả đến hàng phần chục).
	\shortans[]{$3{,}5$}
	\loigiai{
		Điều kiện $ -3<x<5,x\ne 0 $. Khi đó phương trình tương đương
		\begin{eqnarray*}
			\log_2 |x|+\log_2 (5-x)-\log_2(x+3)=0
			&\Leftrightarrow& \log_2 \dfrac{|x|(5-x)}{x+3}=0 \\
			&\Leftrightarrow& \dfrac{|x|(5-x)}{x+3}=1 \\
			&\Leftrightarrow& \heva{&0<x<5\\&x(5-x)=x+3} \vee \heva{&-3<x<0\\&-x(5-x)=x+3} \\
			&\Leftrightarrow& \heva{&0<x<5\\&-x^2+4x-3=0} \vee \heva{&-3<x<0\\&x^2-6x-3=0} \\
			&\Leftrightarrow& x=1 \vee x=3 \vee x=3-2\sqrt{3}.
		\end{eqnarray*}
		Vậy tổng các nghiệm của phương trình là $ 7-2\sqrt{3} \approx 3{,}5$.
	}
\end{ex}

\begin{ex}%[EX chuẩn hóa]%[Trần Chiến]%[1D7B2-1]
	Cho hàm số $f(x) = - \dfrac{1}{3}x^3 + 4x^2 - 7x - 11$. Trong đoạn $[0;2025]$, bất phương trình $f'(x) \geqslant 0$ có tất cả bao nhiêu nghiệm nguyên?
	\shortans[]{$7$}
	\loigiai{
		Ta có $f'(x)=-x^2+8x-7$, do đó $f'(x) \geq 0 \Leftrightarrow -x^2+8x-7 \geq 0 \Leftrightarrow 1\leq x\leq 7$. \\
		Vậy bất phương trình đã cho có $7$ nghiệm nguyên thuộc đoạn $[0;2025]$.	
	}
\end{ex}

\begin{ex}%[EX chuẩn hóa]%[Trần Chiến]%[1H8K6-1]
	Cho hình chóp $S.ABC$, tam giác $ABC$ vuông tại $B$, $SA$ vuông góc với $(ABC)$, $SA = a\sqrt 3 $, $AB = a$. Góc giữa $SB$ và mặt phẳng $(ABC)$ có số đo bằng bao nhiêu độ?
	\shortans[]{$60$}
	\loigiai{\immini{
			Góc giữa $SB$ và mặt phẳng $(ABC)$ là góc $\widehat{SBA}$. \\
			Xét tam giác vuông $SAB$, ta có\\ $\tan \widehat{SBA} = \dfrac{SA}{AB} = \dfrac{a\sqrt{3}}{a} = \sqrt{3}$.\\
			Suy ra $\widehat{SBA} = 60^\circ$. \\
			Vậy góc giữa $SB$ và mặt phẳng $(ABC)$ bằng $60^\circ$.
		}{	
			\begin{tikzpicture}[scale=1, font=\footnotesize, line join=round, line cap=round, >=stealth]
		\path
		(0,0) coordinate (A)
		(3,0) coordinate (C)
		(1,-1) coordinate (B)
		%($(B)!0.5!(C)$) coordinate (I)
		%(intersection of A--M and C--N) coordinate (H)
		($(A)+(0,3)$) coordinate (S);
		\draw (B)--(S)--(A)--(B)--(C)--(S); \draw[dashed] (C)--(A);
		\foreach \i/\g in {S/90, A/-135, B/-90, C/-20}{
			\draw[fill=black] (\i) circle (1.2pt) +(\g:0.25) node{$\i$};}
		
		\draw pic[draw,angle eccentricity=1.7, angle radius=0.4cm]{angle=S--B--A};
		\draw pic[draw,angle radius=2mm]{right angle=A--B--C};
		\end{tikzpicture}	
	}}
\end{ex}

\begin{ex}%[EX chuẩn hóa]%[Trần Chiến]%[1D6K3-2]
	Có bao nhiêu giá trị nguyên dương nhỏ hơn $2022$ của tham số $m$ để hàm số $y=\dfrac{x^2-3 x+5}{\log _{2022}\left(x^2-2x+m^2-4m+5\right)}$ xác định với mọi $x \in \mathbb{R}$?
	\shortans[]{$2018$}
	\loigiai{
		Hàm số đã cho xác với mọi $x \in \mathbb{R}$ khi và chỉ khi
		\begin{eqnarray*}
			& & \heva{& x^2-2x+m^2-4m+5>0\\& x^2-2x+m^2-4m+5\ne 1}, ~\forall x \in \mathbb{R} \Leftrightarrow \heva{& x^2-2x+m^2-4m+5>0, ~\forall x \in \mathbb{R}\\& x^2-2x+m^2-4m+4\ne 0, ~\forall x \in \mathbb{R}}\\
			& \Leftrightarrow & \heva{& \Delta '_1 =1-(m^2-4m+5) <0 \\& \Delta ' _2 = 1-(m^2-4m+4)<0} \Leftrightarrow \heva{& -m^2+4m-3 <1\\& -m^2+4m-3<0 }  \\
			& \Leftrightarrow & -m^2+4m-3<0 \Leftrightarrow m<1 \text{ hoặc } m>3.
		\end{eqnarray*}
		Mà $m \in \mathbb{N}^*$ và $m<2022$ nên $m \in \{4;5;6;\ldots;2021\}$.\\
		Vậy có $2018$ giá trị $m$ cần tìm.
	}
\end{ex}

\begin{ex}%[EX chuẩn hóa]%[Trần Chiến]%[1H8G5-4]
	Cho hình hộp chữ nhật $ ABCD.A'B'C'D' $ có $ AB=AA'=1, AC=2$. Tính khoảng cách giữa hai đường thẳng $ AC' $ và $ CD' $ (kết quả làm tròn đến hàng phần trăm).
	\shortans[]{$0{,}55$}
	\loigiai{
		\immini{Gọi $M,N$ lần lượt là trung điểm của $A'D'$ và $BC$.\\ Dễ thấy $AM \parallel C'N$ nên $A, M, C', N$ đồng phẳng.\\
			Ta có $CD' \parallel MN$ nên $CD' \parallel (AMC'N)$. Suy ra \\
			$\mathrm{d}(CD',AC') = \mathrm{d}(CD',(AMC'N)) = \mathrm{d}(C,(AMC'N))$.\\
			Kẻ $CI \perp AN$ tại $I$ và $CH \perp C'I$ tại $H$, suy ra $\mathrm{d}(C, (AMC'N)) = CH$.\\
			$\triangle CIN \backsim \triangle ABN$ (g-g), suy ra $\dfrac{CI}{AB} = \dfrac{CN}{AN}$ \\
			$\Rightarrow CI = \dfrac{AB \cdot CN}{AN}$.
		}{
		\begin{tikzpicture}[scale=1, font=\footnotesize, line join=round, line cap=round, >=stealth]
			\path
			(0,0) coordinate (A)
			(5,-3) coordinate (C)
			(-1,-3) coordinate (B)
			($(A)+(C)-(B)$) coordinate (D)
			($(A')!1/2!(D')$) coordinate (M)
			($(B)!1/2!(C)$) coordinate (N)
			($(N)!(C)!(A)$) coordinate (I)
			($(C')!(C)!(I)$) coordinate (H)
			($(A)+(0,2)$) coordinate (A')
			($(B)+(0,2)$) coordinate (B')
			($(C)+(0,2)$) coordinate (C')
			($(D)+(0,2)$) coordinate (D');
			\draw[dashed] (D)--(A)--(B) (D)--(C) (M)--(A)--(N)--(C) (D')--(C) (M)--(N) (B)--(A)--(A');
			\draw (B)--(N) (B)--(B') (C)--(C') (M)--(C')--(N)--(I)--(C)--(H) (C')--(I) (D')--(D)--(C) (D')--(C) (A')--(B')--(C')--(D')--cycle;
			\foreach \t/\g in {A/180,B/-90,C/0,D/-80,A'/180,B'/90,C'/0,D'/80,M/90,N/-90,I/-90,H/180}{
				\draw[fill=black] (\t) circle (1pt) node[shift={(\g:7pt)},font=\scriptsize]{$ \t $};
			}
			\draw pic[draw,angle radius=2mm]{right angle=N--I--C}; 
			\draw pic[draw,angle radius=2mm]{right angle=I--H--C}; 
\end{tikzpicture}}
		$BC = \sqrt{AC^2-AB^2} = a\sqrt{3}$; $CN = BN = \dfrac{BC}{2}=\dfrac{a\sqrt{3}}{2}$; $AN = \sqrt{AB^2+BN^2} = \dfrac{a\sqrt{7}}{2}$; $CI = \dfrac{AB \cdot CN}{AN} = \dfrac{\sqrt{21}}{7}$.\\
		Tam giác $C'CI$ vuông tại $C$ có đường cao $CH$, suy ra $CH = \dfrac{CC' \cdot CI}{\sqrt{CC'^2 +CI^2}} = \dfrac{a\sqrt{30}}{10}$.
	}
\end{ex}

\Closesolutionfile{ans}
\indapan{7}{ans/ans-cau-KQ-1}